% =====================================================================
%  Chapter 9 — The Long Tail
% =====================================================================

\chapter{The Long Tail}
\label{ch:longtail}

Five candidate classes have been treated in detail. They are not the only
ones, and the ones that remain are not marginal. This chapter works through
the rest of the space, with the same question in each case: what would have
to be true for this to be dark matter, and what is the cheapest way to find
out whether it is?

\section{Sterile neutrinos}A right-handed neutrino that is not part of the Standard
Model is the simplest extension of the particle content that nobody has
excluded, and the minimal-mass version --- a few eV --- was the original
candidate for the warm dark matter that would solve the small-scale problems.

It has a clear and unfavourable status. The oscillations and the absence of
neutrino-less double beta decay set strong constraints, and the same
production mechanism that populates the neutrino masses in the Standard Model
populates the extra species and produces too much of it at the observed
mass. A sterile neutrino at the eV scale is now strongly constrained by the
combination of cosmology and laboratory data, and the same mechanism that
fails for the eV case also fails for the keV case, because it produces the
wrong spectrum.

There is one window worth watching: keV-scale sterile neutrinos produced
non-thermally, with a different production history. That is a different
model, and it is the kind of model that the X-ray and UV-line searches of
Chapter~\ref{ch:indirect} are designed to find.

\section{SIMP dark matter}In the Simplify Model of Dark Matter, a strong
self-interaction sets the dark matter's effective viscosity, which freezes
out the annihilation rate at a value related to the cross section. The
density is then a prediction of the same interaction that produces the
small-scale structure.

This is a genuinely elegant mechanism, and its testable consequences are
strong. A cross section of order $10\ \cmtwo/\mathrm{g}$ produces cores in
dwarf galaxies, which is why the class has been proposed as a solution to
the cusp-core problem. It also predicts a diversity in the inner density
profiles of dwarf galaxies, which is directly measurable, and a velocity
offset between the stars and the gas in the most ancient stars, which is a
clean early-universe fossil.

The problems are production-related. The SIMP mechanism needs a
three-body annihilation process, which requires a mass scale for the
dark sector well above the kinetic energy of the particles at freeze-out,
and the natural way to achieve that --- a larger cross section in some
channel -- reintroduces the direct-detection and indirect-detection
constraints that the class was designed to avoid. A SIMP with the right
cross section has been shown to be under pressure from the same kind of
constraints.

\section{Freeze-in and the feeble interactions}In the freeze-in scenario, the dark
matter never reaches thermal equilibrium. It is produced very slowly by a
coupling to the Standard Model, and the final abundance is proportional to
that coupling times the expansion rate.

This makes the density a prediction of a coupling that is small by
construction, which resolves the WIMP's miracle problem in the cleanest way
available: the density is not explained, it is parameterised, and the
parameter can be anything. The price is that there is no lower bound on the
interaction strength, and therefore no reason for any particular detection
experiment to be the right one.

Freeze-in does, however, make a generic and useful prediction: the dark
matter is very weakly coupled to the Standard Model, which means it is
essentially undetectable in the laboratory except through its gravitational
effects and through its production in the early universe. If the field
ends up here, the discovery will come from cosmology or from the
astrophysical small-scale programme, and the field's laboratory programme
will have produced null results that are correct and uninformative.

\section{Dark sectors and portals}A hidden sector containing its own fermions and
bosons, coupled to the Standard Model through a portal, is the general
framework into which most of the remaining particle-physics parameter space
falls. The mediator might be a gauge boson, a scalar, or a pseudoscalar.

Two facts about the framework are worth stating. First, it is not a
prediction; it is a category. Second, the portals are already strongly
constrained by the laboratory experiments that were designed to find them
without knowing their mass --- beam-dump, fixed-target, resonance and
precision measurements at the LHC and at dedicated facilities --- so the
surviving portal parameter space is narrow in the mass range where those
experiments are sensitive.

The consequence for this book is a decision about where to put effort. The
dark sector hypothesis is best tested not by searching for the dark matter
directly but by searching for the mediator, which is a collider and beam-dump
problem.

\section{Superheavy dark matter}Dark matter with a mass far above the electroweak scale
can be produced in the early universe by gravitational production or by
thermal production of a massive state that later decays. Such particles have
far too little energy density in the present universe to be a thermal
relic at that mass, so the production mechanism must be non-equilibrium.

The observational windows are unusual. Superheavy dark matter would
gravitationally focus light into extremely bright arcs, and it would
produce a distinctive modulation of the CMB's anisotropies on small scales.
It is also directly accessible to a beam dump: a beam of particles hitting
a foil would produce missing-energy events from the invisible recoiling
particle.

The 2023 BeamDumpBeam experiment at CERN reported a background measurement
rather than a discovery, and the analysis of its first data release is the
clearest statement of what is excluded \idx{Superheavy dark matter}. The
class has one attractive property: the production mechanism for a
superheavy relic is essentially unavoidable in many cosmological scenarios,
which means the abundance is a prediction rather than a parameter.

\section{How the tail should be judged}The classes in this chapter share a
property: none of them can be excluded by the experiments that the field has
been running, and none of them can be confirmed by them either. That is not
a defect of the experiments; it is a statement about where the interesting
parameter space lies.

The right way to think about the tail is as the part of the space where the
laboratory programme has already run out of leverage. The experiments that
remain relevant for these classes are the ones that look for the mediator
rather than the dark matter, the astronomical ones, and the cosmological
ones. That is a portfolio argument in its own right, and it is developed in
Chapter~\ref{ch:portfolio}.

\begin{keytakeaways}
\item Sterile neutrinos, SIMP dark matter, freeze-in, dark sectors and
superheavy relics each fail in a different way, and each fails in a way
that a different experiment would expose.
\item Freeze-in explains the density by parameterising it rather than by
predicting it, which has the consequence that the candidate is essentially
undetectable in the laboratory.
\item The dark sector hypothesis is best tested by searching for the
mediator, which is a collider and beam-dump problem rather than a
direct-detection problem.
\item None of the classes in this chapter is excluded or confirmed by the
laboratory experiments the field has been running.
\end{keytakeaways}
